% Generated draft: requires mathematical and source review.
\documentclass[11pt]{article}
\usepackage[margin=25mm]{geometry}
\usepackage{fontspec}
\setmainfont{Latin Modern Roman}
\usepackage{amsmath,amssymb,mathtools,unicode-math}
\setmathfont{Latin Modern Math}
\usepackage{enumitem,needspace,xurl,hyperref}
\setlength{\parindent}{0pt}
\setlength{\parskip}{5pt}
\setlength{\emergencystretch}{3em}
\newcommand{\N}{\mathbb{N}}
\newcommand{\Z}{\mathbb{Z}}
\newcommand{\Q}{\mathbb{Q}}
\newcommand{\R}{\mathbb{R}}
\newcommand{\C}{\mathbb{C}}
\newcommand{\F}{\mathbb{F}}
\newcommand{\A}{\mathbb{A}}
\newcommand{\PP}{\mathbb{P}}
\newcommand{\E}{\mathbb{E}}
\newcommand{\eps}{\varepsilon}
\newcommand{\dd}{\,\mathrm{d}}
\begin{document}
\section*{UnsolvedMath ID: 30005157}
\textbf{Source problem code:} catalogue entry
\subsection{Definitions and mathematical statement}
The two-dimensional moment curve is \(\gamma(t)=(t,t^2)\). Write \(J\subset\R\) for the parameter interval intended by the source. With the conventional Fourier normalization, for integrable \(f:J\to\C\) define
\[
Ef(x)=\int_J f(t)e^{2\pi i(x_1t+x_2t^2)}\,\dd t,
\qquad
E_If(x)=\int_I f(t)e^{2\pi i(x_1t+x_2t^2)}\,\dd t,
\quad x=(x_1,x_2)\in\R^2.
\]
This normalization makes the notation explicit; the supplied excerpt does not specify the source's normalization or integration measure. For an admissible scale \(\delta\), let \(\mathcal I(\delta)\) be the source's partition of \(J\) into intervals of length \(\delta\), and let \(B_{\delta^{-2}}\subset\R^2\) be the source's ball of radius \(\delta^{-2}\). For a specified nonnegative measurable weight \(w_B\), write
\[
\|g\|_{L^p(w_B)}=\left(\int_{\R^2}|g(x)|^p w_B(x)\,\dd x\right)^{1/p},
\qquad
\|g\|_{L^p(B)}=\left(\int_B|g(x)|^p\,\dd x\right)^{1/p}.
\]
The weight \(w_B\) is intended to be the particular weight prescribed by the source, not an arbitrary weight.

The supplied question asks whether there exist an exponent \(p\in(2,4)\) and a finite constant \(C_p\), independent of \(\delta\) and \(f\), such that, for every admissible scale \(\delta\) and every function \(f\) in the intended function class,
\[
\|Ef\|_{L^p(B_{\delta^{-2}})}
\leq C_p
\left\|\left(\sum_{I\in\mathcal I(\delta)}|E_If|^2\right)^{1/2}\right\|_{L^p(w_{B_{\delta^{-2}}})}.
\]
The quantifier on \(p\) is existential: the question requests at least one exponent strictly between \(2\) and \(4\). The record describes this as removing logarithmic losses from estimates associated with equation (5) of the report. The complete equation (5) and the precise logarithmic factors are not supplied, so that description cannot independently establish an equivalence between fully specified estimates.

The displayed question remains partially specified. The supplied record does not recover \(J\), the intended function class, the exact admissible scales, endpoint conventions for the partition, the formula and normalization of \(w_B\), or the quantifier over ball centers. These are unresolved hypotheses, not choices made by this draft. In particular, the interval \([0,1]\) appearing in the following background proposition does not by itself specify the extension operator's parameter interval.

For context, the surrounding excerpt attributes a \(2n\)-orthogonality argument to C\'ordoba and Fefferman and says that the proof of its Theorem 1 reduces to Proposition 2, cited there as [2]. Reading the corrupted power notation in the supplied extraction as powers of coordinates, that proposition states: for every \(n\in\N\), there exists \(C_n\geq1\) such that, whenever \(0<\delta<1\) and \(x,y\in[0,1]^n\) satisfy
\[
\left|\sum_{r=1}^n x_r^k-\sum_{r=1}^n y_r^k\right|\leq\delta^n
\qquad(1\leq k\leq n),
\]
there exists a permutation \(\sigma\) of \(\{1,\ldots,n\}\) for which
\[
|x_r-y_{\sigma(r)}|\leq C_n\delta
\qquad(1\leq r\leq n).
\]
This is supplied background about almost solutions to a Vinogradov-type system, not the open question to be established. Its typography requires checking against the report; the statement of Theorem 1 and the bibliographic identity of [2] are absent.
\subsection{Short English statement}
The supplied problem asks whether the Fourier extension operator for the planar moment curve admits the displayed square-function estimate for at least one \(2<p<4\), with a constant uniform in the scale \(\delta\). The parameter interval, weight, function class, scale conventions, and ball-center quantifiers require recovery before this becomes a fully specified statement. The supplied open status and logarithmic-loss assessment are dated curation, not independently verified conclusions.
\subsection{Sources}
\begin{itemize}
\item Local UnsolvedMath registry: supplied complete record for permanent ID \(30005157\), title \emph{Uniform Intermediate-Norm Square-Function Estimates}, exact source code \texttt{catalogue entry}, in the Oberwolfach Reports Open Problems collection. No individual proposer is supplied; the proposed year is recorded as 2022.
\item Supplied report citation: \emph{Real Analysis, Harmonic Analysis and Applications} (2023), Oberwolfach Reports, pp.\ 1743--1803, DOI \texttt{10.4171/owr/2022/31}. \url{https://doi.org/10.4171/owr/2022/31}. The extracted passage contains the page marker 1767.
\item Supplied literature source, labeled ``OWR report'': \url{https://publications.mfo.de/bitstream/handle/mfo/4013/OWR_2022_31.pdf?sequence=4}.
\item Literature assessment dated exactly \texttt{2026-08-22}: supplied dated curation, not independently verified. It labels the problem open, reports logarithmic-loss alternatives, and says that no source located in that assessment removes the loss. The background explicitly requests independent verification of time-sensitive status and cited claims.
\item The registry supplies the label \texttt{corrected\_verified} and reports that the official report was checked and its notation repaired. These are supplied verification claims, not verification performed for this draft. The surrounding excerpt attributes the orthogonality argument to C\'ordoba and Fefferman and labels its contextual proposition ``Proposition 2 ([2])''; no bibliographic entry for [2] is supplied.
\end{itemize}
\end{document}
